\documentclass[10pt]{beamer}
\usetheme[numbering=fraction, background=light]{metropolis}

\usepackage{graphicx}
\usepackage{booktabs}
%\usepackage{../Self_Style}
\usepackage{amsmath, amsthm, amssymb}

\usepackage[normalem]{ulem} %for strike through

% No fontspec, stick with default font

\title{Intro to Peano Systems}
\subtitle{Y'all love Set Theory, don't you?}
\author{Mysterious Category Theory Enthusiast}
%\institute{UCSB, CCS Mathematics}
\date{\today}

\setbeamertemplate{theorems}[ams style]
% \newtheorem{theorem}{Theorem}
% \newtheorem{corollary}[Theorem]{Corollary}
%\newtheorem{definition}[Theorem]{Definition}
% \newtheorem{example}[Theorem]{Example}
% \newtheorem{lemma}[Theorem]{Lemma}
\newtheorem{proposition}[theorem]{Proposition}
\newtheorem{remark}[theorem]{Remark}
\newtheorem{question}{Question}
\newtheorem{property}{Property}

\begin{document}
\maketitle

\begin{frame}{What is Inductive Set?}
    \begin{definition}
        Given a set $x$, its \emph{Successor} $x^+:= x\cup \{x\}$. 
    \end{definition}

    \hfil

    \begin{definition}
        A set $I$ is an \emph{Inductive Set}, if:
        \begin{itemize}
            \item[1.] $\emptyset\in I$.
            \item[2.] For any $x\in I$, $x^+\in I$.
        \end{itemize}
    \end{definition}

    \hfil

    \begin{example}
        Define $0:= \emptyset$, and $(n+1):= n^+ = n\cup \{n\}$. This realizes $\mathbb{N}$ as an inductive set.
    \end{example}
\end{frame}

%\begin{frame}{Morphisms of Inductive Sets}
%    Let $I,J$ be two inductive sets.

%    \begin{definition}
%        A \emph{morphism of inductive sets} $f:I\rightarrow J$, is a set function, such that 
%        \[f(x)^+ = f(x^+)\]
%        for any set $x\in I$.
%    \end{definition}
%\end{frame}

\begin{frame}{Properties of Inductive Sets}
    \begin{proposition}
        Unions / Intersections of inductive sets is inductive set.
    \end{proposition}

    \hfil

    \begin{proof}
        Fill in the details on the blackboard :)
    \end{proof}
\end{frame}

\begin{frame}{Is there a SMALLEST Inductive Set?}
    Surely there is one, right? Like natural numbers $\mathbb{N}$.

    \hfil

    \sout{\textbf{Answer 1:} Yes, take intersection of all inductive sets.}

    \hfil \sout{Yep, this answer is screwed up by ZFC axioms.}

    \hfil

    \textbf{Answer 1':} Kind of: Take $S$ as a set of (arbitrary) inductive sets, define $\omega:= \bigcap_{I\in S}I$, this is clearly inductive. Question is, is it the smallest?
\end{frame}

\begin{frame}{$\omega$ is the SMALLEST Inductive Set}
    Fix any inductive set $T$, define $S_T:=$ all inductive subsets of $T$. Defin $\omega$ with respect to $S:= S_T$.

    \hfil

    \begin{theorem}
        For each inductive set $I$, $\omega\subseteq I$.
    \end{theorem}

    \hfil

    In particular, $\omega\cong \mathbb{N}$, surprise, surprise...
\end{frame}

\begin{frame}{Proof}
    For any inductive set $B$, consider $\omega\cap B\subseteq \omega\subseteq T$.

    Since $\omega\cap B$ is also inductive, $\omega\cap B\in S_T$, so $\omega\subseteq \omega\cap B$, implying $\omega=\omega\cap B$. 

    Hence, $\omega\subseteq B$. $\square$
\end{frame}

\begin{frame}{Peano System}
    \begin{definition}
        A \emph{Peano System} is a set $N$, an element $e\in N$, and a function $s:N\rightarrow N$, such that:
        \begin{itemize}
            \item $e\notin s(N)$.
            \item If $A\subseteq N$, $e\in A$, and $s(A)\subseteq A$, then $A=N$.
        \end{itemize}
    \end{definition}

    \hfil

    \begin{example}
        The set $\omega$ of all finite successors of $\emptyset$, with successing function $(\_)^+:\omega\rightarrow\omega$ is a Peano System. In particular, $\omega\cong \mathbb{N}$ as Inductive Sets / Peano Systems.

        Moreover, $\omega$ is the \emph{initial} Peano System, surprise, surprise...
    \end{example}
\end{frame}

\end{document}